\selectlanguage{ngerman}
\begin{ledgroupsized}[r]{120mm}
\footnotesize
\pstart
\noindent\textbf{Überlieferung:}
\pend
\end{ledgroupsized}
\begin{ledgroupsized}[r]{114mm}
\footnotesize
\pstart \parindent -6mm
\makebox[6mm][l]{\textit{L}}%
Aufzeichnung:
LH~XXXVII~5 Bl.\;25.
Ein Zettel (etwa 5,5\;x\;10~cm),
Ränder beschnitten;
Papiererhaltungsmaßnahmen;
\href{\Scan{LH_37_05_025wz}}{Harzer Papier}.
Eine Seite auf Bl.\;25\,r\textsuperscript{o}.
Bl.\;25\,v\textsuperscript{o} überliefert lediglich das Diagramm \lbrack\textit{Fig.\:1}\rbrack,
welches keinen erkennbaren Bezug zum Text aufweist.
\pend%
\end{ledgroupsized}
\vspace{5mm}
\begin{ledgroup}
\footnotesize
\pstart
\noindent%
\textbf{Datierungsgründe:}
Die Datierung der Aufzeichnung N.~\ref{60283} beruht auf äußeren Eigenschaften des Text\-trägers %
sowie auf inhaltlichen Betrachtungen.
Siehe die ausführliche Begründung auf S.~\pageref{60283_dat}.
\pend%
\end{ledgroup}%
\selectlanguage{latin}%
\frenchspacing%
\vspace{8mm}%
\normalsize%
\pstart%
\noindent%
\lbrack%
\href{\Scan{LH_37_05_025r}}%
{25\:r\textsuperscript{o}}%
\rbrack\ %
Corpus impingens%
\indexS{corpus impingens}
in penetrabile%
\indexS{corpus penetrabile}
cum difficultate,%
\indexS{difficultas penetrandi}
tanta celeritate ferri potest,%
\indexS{celeritas corporis impingentis}
ut ab eo repercutiatur.%
\indexS{repercussio corporis impingentis}
Nam corpus penetrabile ipsummet%
\indexS{corpus penetrabile}
tanta celeritate moveri potest,%
\indexS{celeritas corporis penetrabilis}
ut aliud solidum sustineat,%
\indexS{corpus solidum}
imo et ut fortius propellat,
quam ipsummet
\edlabel{LH_37_05_025_025r_evbw-1}%
fertur.%
\edtext{}{%
\xxref{LH_37_05_025_025r_evbw-1}{LH_37_05_025_025r_evbw-2}%
\lemma{fertur.}%
\Bfootnote{%
\textit{(1)}~Causa %
\textit{(2)}~Idem autem %
\lbrack...\rbrack\ effectum. \lbrack/\rbrack\ Causa%
~\textit{L}%
}}
Idem autem est
corpus impingere et excipere%
\lbrack,\rbrack\
quantum ad hunc effectum.
\pend%
\pstart%
Causa%
\edlabel{LH_37_05_025_025r_evbw-2}
horum omnium,
quod eo
\edtext{%
major est difficultas penetrandi,%
\indexS{difficultas penetrandi}
quo major est penetrantis celeritas;%
\indexS{celeritas penetrantis}%
\indexS{corpus penetrans}%
}{%
\lemma{major est}%
\Bfootnote{%
\textit{(1)}~celeritas %
\textit{(2)}~difficultas penetrandi, %
\lbrack...\rbrack\ penetrantis celeritas;%
~\textit{L}%
}}
et potest esse impingentis celeritas%
\indexS{corpus impingens}%
\indexS{celeritas corporis impingentis}
tanta,
ut citius ipsum
\edtext{%
corpus excipiens%
\indexS{corpus excipiens}
flectatur et resiliat,
adeoque et repercutiat,
quam ut%
}{%
\lemma{corpus}%
\Bfootnote{%
\textit{(1)}~impingens flectatur et repercutiatur %
\textit{(2)}~excipiens flectatur et resiliat, %
\textit{(a)}~quam ut %
\textit{(b)}~adeoque et %
\textit{(aa)}~repercutiatur, %
\textit{(bb)}~repercutiat, quam ut%
~\textit{L}%
}}
dividi se patiatur.
Nimirum divisio%
\indexS{divisio corporis penetrabilis}
ex eo oritur,
quod totum motui magis resistit,%
\indexS{motus corporis impingentis}
quam elastrum partes connectens.%
\indexS{elastrum connectens partes}
\lbrack%
\href{\Scan{LH_37_05_025v}}%
{25\:v\textsuperscript{o}}%
\rbrack%
\pend%
\vspace{1.25em}%
\pstart%
\noindent%
\lbrack\textit{Ohne erkennbaren Bezug zum Text:}\rbrack\
\pend%
\vspace{0.50em} %
\centerline{%
\hspace{1.5mm}%
\includegraphics[width=0.060\textwidth]{%
gesamttex/edit_VIII-4/images/LH_37_05_025_d_025v.pdf%
}} 
\nobreak\vspace{0.25em}
\centerline{%
\hspace{2.0mm}%
\lbrack\textit{Fig.\:1}\rbrack%
}
\newpage%
\vspace*{1.0em}%
\pstart%
\noindent%
\centering%
*     *     *
\pend%
\vspace{2.0em}%
\selectlanguage{latin}%
\frenchspacing%
\begin{ledgroup}%
\footnotesize%
\pstart%
\noindent%
\textls{Ausführliche Datierungsgründe:}%
\pagelabel{60283_dat}
Die titellose und undatierte Aufzeichnung N.\;\ref{60283} ist mit der These befasst, dass ein weicherer Körper einen eindringenden härteren zurückzustoßen vermöge, wenn die Geschwindigkeit des Zusammenpralls hinreichend hoch sei.
Die These wird durch die allgmeine Annahme begründet, dass umso mehr Widerstand einem Körper beim Eindringen in einen anderen Körper widerfahre, je höher die relative Geschwindigkeit sei.
Die für N.\;\ref{60283} vorgeschlagene Datierung lässt sich folgendermaßen begründen:
\pend%
\pstart%
Der Träger der vorliegenden Aufzeichnung weist zwar kein Wasserzeichen auf, kann aber anhand des gegitterten \href{\Scan{LH_37_05_025wz}}{Papiermusters} einer Harzer\indexO{Harz} Papiermühle zugeordnet werden. %
Papier dieser Art ist in Leibnizens Nachlass für den Zeitraum 1680--1684 belegt.
Demzufolge erweist sich als wahrscheinlich, dass auch N.\;\ref{60283} in diesem Zeitraum entstand.
Der in N.\;\ref{60283} ausgrdruckte Gedanke findet sich tatsächlich in Texten der frühen 1680er Jahre wieder, so etwa in dem zwischen August 1681 und Ende 1682 abgefassten Entwurf \textit{De soni generatione, propagatione et expressione in organo, mechanice explicatis} (\textit{LSB} VIII,\;3 N.\;12\textsubscript{1}, S.\;96.2--4; 98.10--14)\cite{02527} und in dem anschließenden Konzept \textit{L\textsuperscript{1}} zum Entwurf \textit{Cogitationes novae quomodo formetur sonus et per aerem propagetur atque in organo auditus exprimatur} (\textit{LSB} VIII,\;3 N.\;12\textsubscript{3}, S.\;118.11--119.2),\cite{02529} welches spätestens im Winter 1682/1683 angefertigt wurde.
Der Gedanke lässt sich aber auch in älteren Texten nachweisen, die auf Leibnizens frühere Auseinandersetzung mit Honoré Fabris \textit{Physica, id est scientia rerum corporearum} (4 Bde, Lyon 1669-1671)\cite{00044} zurückgehen; vgl. etwa den wohl zwischen Ende 1674 und Ende 1677\;(?) entstandene Entwurf \textit{De motu} (\textit{LSB} VIII,\;3 N.\;4, S.\;14.4--8)\cite{02598} sowie die auf die Zeitspanne Frühjahr 1679 bis Winter 1680/1681\;(?) datierten Aufzeichnung \textit{LSB} VIII,\;3 N.\;7 (Randbemerkung zu S.\;25.4--7).\cite{02599}
Die Aufzeichnung N.\;\ref{60283} könnte somit auch etwas früher als 1680 entstanden sein. 
Auch eine etwa spätere Datierung als 1684 ist möglich.
\pend%
\end{ledgroup}%
\selectlanguage{latin}%
\frenchspacing%
\normalsize%
\newpage%
